% NewsMead core manuscript inserts
% Generated from the implementation-alignment documentation and verified
% against ReadingStateInferencer, WindowedStabilityEstimator,
% AdaptiveScaffoldController, and StudyConfig.
%
% This file is an insertion aid, not a standalone LaTeX document.

\subsection{Session Reading State Index}

The Session Reading State Index (Session RSI) is a cumulative measure of
detected reading effort across the current article-reading session. It combines
regression, dwell, and fixation components on a 0--100 scale using weights of
0.40, 0.35, and 0.25, respectively. Higher values indicate that the gaze-derived
reading pattern contains more behaviors associated with rereading, hesitation,
or effort. Session RSI is a descriptive research outcome; it is not a
comprehension score and does not directly select visual-scaffolding levels.

\begin{equation}
  \mathrm{Session\ RSI}=100(0.40R+0.35D+0.25F),
\end{equation}

where \(R\) is the confirmed regression count divided by the number of lines
reached, bounded to \([0,1]\); \(D\) is accumulated dwell time divided by
elapsed session time, bounded to \([0,1]\); and \(F\) is fixation rate divided
by a reference rate of 120 fixations per minute, bounded to \([0,1]\).

Regression receives the largest engineering weight because confirmed backward
line movement is treated as the strongest of the three available indicators of
rereading. Dwell receives the second-largest weight, and fixation rate the
smallest, because both can also reflect ordinary careful reading or individual
reading style. These weights encode the prototype's current design priorities;
they must not be presented as population-optimal coefficients and should be
examined during pilot validation.

A cumulative measure is unsuitable for immediate adaptation. Difficulty or
tracking noise early in an article may keep Session RSI elevated after the
reader recovers, while a new difficulty episode late in a long stable session
may be diluted by earlier history. Session RSI is also not normalized against
each participant's normal reading pattern. It is therefore retained for
end-of-session description and comparison of articles or experimental
conditions rather than real-time scaffold selection.

\subsubsection{Adaptive Instability Index}

The Adaptive Instability Index (AII) is a recent, participant-relative measure
used for real-time scaffold selection. It is computed from regression rate,
dwell-time proportion, and fixation rate within a rolling 10-second window.
Every 500~ms, after at least two seconds of recent observations are available,
each metric is standardized against the participant's baseline. Only positive
deviations in the direction associated with increased effort are retained.

For metric \(i\), the positive standardized deviation is

\begin{equation}
 z_i^{+}=\max\left(0,
 \frac{x_{i,\mathrm{recent}}-\mu_{i,\mathrm{baseline}}}
 {\max(\sigma_{i,\mathrm{baseline}},\sigma_{i,\mathrm{floor}})}
 \right).
\end{equation}

The raw index is

\begin{equation}
 \mathrm{AII}_{\mathrm{raw}}=
 \operatorname{clamp}_{[0,4]}
 \left(0.40z_R^{+}+0.35z_D^{+}+0.25z_F^{+}\right).
\end{equation}

The implemented standard-deviation floors are 1.0 regression per minute, 0.05
dwell-time proportion, and 5.0 fixations per minute. The raw value is
exponentially smoothed with a two-second time constant. For update interval
\(\Delta t\), \(\alpha=1-e^{-\Delta t/2s}\), and

\begin{equation}
 \mathrm{AII}_t=\mathrm{AII}_{t-1}
 +\alpha(\mathrm{AII}_{\mathrm{raw},t}-\mathrm{AII}_{t-1}).
\end{equation}

The AII is reported in weighted z-derived index units. It must not be
interpreted as an exact number of standard deviations or assigned a
standard-normal percentile because its components are truncated at zero,
weighted, calculated from overlapping windows, capped, and smoothed.

The current prototype automatically estimates a fresh baseline during the first
20 seconds of each article. This is an engineering fallback and is not the
intended main-study baseline protocol. For the main study, each participant will
first read standardized passages under the intended device position, lighting,
font, layout, and posture. The resulting mean and standard deviation of each
metric will be persisted and reused across the participant's experimental
articles. Baseline duration and test-retest reliability will be checked during
the pilot before the profile is used in the main evaluation.

\subsection{Adaptation and Graded Visual Scaffolding}

Visual-scaffolding levels are selected using the smoothed Adaptive Instability
Index rather than the cumulative Session RSI. The controller applies only one
intervention at a time and escalates or withdraws support conservatively to
avoid rapid visual oscillation.

\begin{table}[ht]
\centering
\caption{Adaptive Instability Index mapping used by the current prototype.}
\begin{tabular}{p{0.36\linewidth}p{0.53\linewidth}}
\hline
\textbf{Adaptive Instability Index} & \textbf{Desired state}\\
\hline
Less than 1.0 & No scaffold\\
1.0 to less than 1.5 & Level 1: word highlighting\\
1.5 to less than 2.0 & Level 2: line emphasis\\
2.0 to less than 2.5 & Level 3: five-line focus window\\
At least 2.5 with loss evidence & Level 4: re-entry support\\
At least 2.5 without loss evidence & Level 3: five-line focus window\\
\hline
\end{tabular}
\end{table}

Adaptation remains disabled until baseline collection is complete. At least 65
percent of gaze samples in the recent window must land on visible article text.
This value is a coverage gate, not proof that the correct line or word was
identified. Elevated behavior must persist for 1.5 seconds before escalation,
while recovery must persist for 3.5 seconds before withdrawal. Both escalation
and withdrawal proceed one level at a time. Low gaze coverage clears the active
intervention instead of allowing possible tracking loss to be treated as reading
difficulty.

Level 4 requires evidence beyond an elevated AII. In the current operational
definition, gaze must remain off the visible article text for at least 800~ms
and then return at least two lines away from the last stable line. The resulting
re-entry evidence remains eligible for 10 seconds. A high index or a regression
alone does not activate Level 4.

The thresholds 1.0, 1.5, 2.0, and 2.5; the 65 percent coverage gate; the 1.5-
and 3.5-second persistence intervals; and the current loss-of-position rule are
provisional engineering parameters. The final parameters will be tuned and
frozen using pilot data with independently labeled stable, difficulty,
tracking-error, loss-of-position, and recovery episodes.

\paragraph{Alternative considered.}
Direct mapping of cumulative Session RSI to scaffold levels was considered but
not adopted. Its historical accumulation can preserve early difficulty after
recovery, dilute a late difficulty episode after prolonged stable reading, and
apply non-personalized absolute cutoffs to readers with different normal reading
patterns. Direct control would also require a new set of empirically validated
0--100 thresholds; the current AII thresholds cannot be converted directly to
Session RSI ranges. A high Session RSI would not replace the independent
loss-of-position evidence required for Level 4.
